% !TEX encoding = UTF-8 Unicode
% !TEX root = thesis.tex

\chapter{Results and Discussion}
\label{ch:results}

This chapter reports the results of the four analysis layers and the external
validation, and discusses what they mean. Unless stated otherwise, all numbers
are computed on the \texttt{timeline\_medium} timeline (30 phases, 1\,800\,s,
eight drifts, forty anomalies) with the protocols of
Chapter~\ref{ch:method}. Each section maps to one research question:
anomaly detection (RQ1), drift detection (RQ2), drift-aware maintenance (RQ3),
and the configuration--performance link with reconfiguration (RQ4); a final
section reports how well the findings generalise.

\section{Anomaly detection (RQ1)}
\label{sec:res-anomaly}

Table~\ref{tab:anomaly-type} reports the per-type precision--recall AUC (PR-AUC)
of the five detectors, and Table~\ref{tab:anomaly-overall} the pooled PR-AUC over
the whole timeline with 95\% bootstrap intervals. The pattern is fairly
consistent: the two anomalies that leave a clear, structured signature are
detected reasonably well, while the two subtle ones sit close to the noise floor.

\begin{table}[ht]
\centering
\caption{Per-type anomaly PR-AUC at the 5\,s window (best per row in bold; own elaboration).}
\label{tab:anomaly-type}
\begin{tabular}{@{}lrrrrr@{}}
\toprule
Anomaly type & IsoForest & LOF & OCSVM & PCA-AE & MLP-AE \\
\midrule
Measurement glitch  & 0.635 & 0.842 & \textbf{0.858} & 0.784 & 0.756 \\
RLF burst           & 0.616 & 0.703 & 0.685 & \textbf{0.751} & 0.655 \\
Interference spike  & 0.069 & \textbf{0.152} & 0.099 & 0.111 & 0.078 \\
Slow degradation    & 0.045 & 0.038 & 0.045 & \textbf{0.048} & 0.036 \\
\bottomrule
\end{tabular}
\end{table}

\begin{table}[ht]
\centering
\caption{Pooled anomaly PR-AUC over the whole timeline, with 95\% bootstrap CI (own elaboration).}
\label{tab:anomaly-overall}
\begin{tabular}{@{}lr@{}}
\toprule
Detector & Pooled PR-AUC [95\% CI] \\
\midrule
IsoForest   & 0.425 [0.391, 0.460] \\
PCA-AE      & 0.262 [0.237, 0.287] \\
MLP-AE      & 0.236 [0.212, 0.262] \\
LOF         & 0.157 [0.149, 0.167] \\
OneClassSVM & 0.154 [0.145, 0.164] \\
\bottomrule
\end{tabular}
\end{table}

Measurement glitch and RLF burst are recovered with PR-AUC in the $0.6$--$0.86$
range, because both leave a structured fingerprint across several features (a
frozen RSRP, or a correlated SINR/RSRQ collapse with extra RLF events).
Interference spike is much harder (best PR-AUC $0.152$, LOF): it is cell-targeted
and short, so only a fraction of the fleet sees it in any window. Even so, these
values must be read against the imbalance of the task
(Section~\ref{sec:anomaly-method}): with a positive prevalence of only about
$2.4\%$ for interference spike, the random-classifier PR-AUC baseline is
$\approx 0.024$, so LOF's $0.152$ is roughly six times the chance level, not a
non-detection. Slow degradation is the only type that stays essentially at
that floor (PR-AUC $\approx 0.04$ against a $\approx 0.037$ baseline) because a
slow ramp on a single UE barely moves a fleet-aggregated window. This separation between
``easy'' and ``subtle'' anomalies is the main observation for RQ1, and is part of
what motivates the adaptive layer of Section~\ref{sec:res-adaptive}.

No single detector performs best everywhere, which is worth noting. By per-type
mean PR-AUC the strongest detector is LOF, whereas by pooled PR-AUC it is
Isolation Forest ($0.425$); the two views differ because the pooled metric is
dominated by the large number of normal windows and by the easy anomalies. A
reasonable reading is that detector choice should follow the anomaly profile of
interest, not a single headline number.

\paragraph{Robustness checks.}
The window-size sweep ($2$--$30$\,s) shifts the per-type numbers but not the
overall picture: the shortest window helps the sharp, localised anomalies
(measurement-glitch PR-AUC rises to $0.919$ for One-Class SVM and interference
spike to $0.195$ for LOF at $2$\,s), while the longest window helps the
anomalies that need time to accumulate evidence (RLF-burst rises to $0.845$ for
PCA-AE at $30$\,s, and slow degradation roughly doubles, though it stays near
the noise floor). The drop-one-family ablation shows that
the RSRQ feature family is the most important by a wide margin: removing it costs
up to $0.33$ PR-AUC on measurement glitch, whereas removing the event-count or
derived-rate families changes little. Figure~\ref{fig:anomaly-degradation} shows
the per-phase false- and true-positive rates with drift phases shaded, making
visible how a static detector drifts off its operating point as the environment
changes.

\begin{figure}[ht]
\centering\includegraphics[width=.85\textwidth]{drift_degradation}
\caption{Per-phase false-/true-positive rate of the static anomaly monitor, drift phases shaded (own elaboration).}
\label{fig:anomaly-degradation}
\end{figure}

\section{Drift detection (RQ2)}
\label{sec:res-drift}

Every one of the eight injected drifts is detected by every one of the seven
detectors (instance-level miss rate $0$). What separates the detectors is the
price: how late the alarm comes, and how often it fires on clean data.
Table~\ref{tab:drift} gives the median detection latency per drift type together
with the mean baseline false-alarm rate, and Figure~\ref{fig:drift-heatmap}
visualises the same trade-off.

\begin{table}[ht]
\centering
\caption{Median drift-detection latency (s) per drift type, and mean baseline false-alarm rate (per second) on clean phases (own elaboration).}
\label{tab:drift}
\begin{tabular}{@{}lrrrrr@{}}
\toprule
Detector & Channel & Mobility & Reconfig & Traffic & FPR\,(/s) \\
\midrule
Energy-batch & 3.0  & 3.0  & 3.0  & 3.0  & 0.115 \\
MMD-batch    & 3.0  & 3.0  & 3.0  & 3.0  & 0.117 \\
DDM          & 12.0 & 2.0  & 3.0  & 15.5 & 0.005 \\
Page--Hinkley & 17.0 & 10.5 & 6.0  & 11.5 & 0.011 \\
EDDM         & 25.5 & 3.0  & 7.5  & 5.0  & 0.006 \\
ADWIN        & 23.0 & 19.0 & 13.0 & 29.0 & 0.005 \\
KSWIN        & 51.5 & 20.0 & 20.0 & 24.5 & 0.008 \\
\bottomrule
\end{tabular}
\end{table}

\begin{figure}[ht]
\centering\includegraphics[width=.9\textwidth]{drift_detection_heatmap}
\caption{Drift-detection latency and baseline false-alarm rate per detector and stream (own elaboration).}
\label{fig:drift-heatmap}
\end{figure}

The results point to a latency--false-alarm trade-off. The two batch two-sample
tests, Energy-batch and MMD-batch, react fastest: they flag every drift type in
about $3$\,s, but at the price of a high baseline false-alarm rate
($\approx 0.12$ per second). The sequential and error-rate detectors are the
opposite: ADWIN and DDM keep the baseline false-alarm rate around $0.005$ per
second but take longer, and DDM is notably quick on the abrupt mobility and
reconfiguration drifts ($2$--$3$\,s) where a sharp mean change is present. By the
combined latency criterion (with misses penalised) the overall winner is
Energy-batch, but in an operational setting the low-false-alarm detectors would
be preferable; the right choice depends on whether the cost of a late alarm or a
false alarm is higher. This trade-off, not the crowning of a single best detector, is the
main observation for RQ2.

\section{Drift-aware adaptive monitoring (RQ3)}
\label{sec:res-adaptive}

The PCA-reconstruction detector (PCA-AE) is used as the base monitor for the
adaptive experiments rather than the pooled-best Isolation Forest. The reason is
that the adaptive loop repeatedly refits the detector under drift, so the
selection criterion here is not the highest static accuracy but suitability for
frequent, stable retraining: PCA-AE yields a single continuous
reconstruction-error score that is cheap to refit and behaves consistently across
refits, which makes the quantile-based filtering step (dropping the most
suspicious windows before a refit) well defined. A tree-based score such as
Isolation Forest re-grows its ensemble at every refit and shifts its operating
threshold, which is less suited to the self-supervised filtering used below.

Table~\ref{tab:adaptive} compares the four retraining strategies on the PCA-AE
monitor, and Figure~\ref{fig:adaptive} plots PR-AUC over time. The four
strategies share the same base detector and the same trigger, yet their outcomes
diverge sharply: what appears to decide the result is what is allowed into the
retraining pool.

\begin{table}[ht]
\centering
\caption{Retraining strategies on the PCA-AE monitor (own elaboration).}
\label{tab:adaptive}
\begin{tabular}{@{}lrrrr@{}}
\toprule
Strategy & Overall PR-AUC & Drift-phase PR-AUC & Refits & Train samples \\
\midrule
Static                  & 0.264 & 0.057 & 1 & 2\,016 \\
Periodic                & 0.115 & 0.036 & 9 & 14\,088 \\
Drift-triggered naive   & 0.116 & 0.037 & 8 & 12\,048 \\
Drift-triggered filtered & 0.248 & \textbf{0.187} & 8 & 10\,040 \\
\bottomrule
\end{tabular}
\end{table}

\begin{figure}[ht]
\centering\includegraphics[width=.85\textwidth]{adaptive_pr_auc_over_time}
\caption{Sliding PR-AUC over time for the four strategies, with the cost-versus-gain summary (own elaboration).}
\label{fig:adaptive}
\end{figure}

Naive retraining tends to backfire here. Both periodic and
drift-triggered-naive retraining roughly halve the overall PR-AUC (from $0.264$
to $\approx 0.12$) and barely move the drift-phase PR-AUC. A likely reason is
self-poisoning: because the retraining pool is taken without labels, a refit
triggered just after a drift absorbs the post-drift anomalies into its notion of
``normal'', so the model learns to ignore much of what it should flag. The
self-supervised filter largely addresses this. The drift-triggered-filtered
strategy keeps the overall PR-AUC close to the static baseline ($0.248$ vs
$0.264$) and more than triples the drift-phase PR-AUC ($0.187$ vs $0.057$), a
$+0.13$ improvement that meets the acceptance target of $+0.10$. It also does so
at lower cost than the naive variants (fewer training samples), because the
filter discards the most suspicious windows before refitting. The observation for
RQ3 is therefore a conditional one: in these experiments drift-aware retraining
helps mainly when the retraining data is first cleaned of likely anomalies, while
an unfiltered refit performs worse than not retraining at all.

\section{Configuration--performance and reconfiguration (RQ4)}
\label{sec:res-config}

\paragraph{Surrogate accuracy.}
Table~\ref{tab:surrogate-point} reports the cross-validated accuracy of the
HistGB surrogate on the three mobility KPIs. HOSR and RLF rate are predicted
accurately ($R^2 = 0.96$ and $0.95$; HOSR MAE $0.043$), and ping-pong rate is
harder but still useful ($R^2 = 0.72$). Because the cross-validation is grouped
by configuration, these scores describe performance on configurations the model
has never seen, not memorised training points.

\begin{table}[ht]
\centering
\caption{Surrogate point accuracy, grouped 5-fold CV (own elaboration).}
\label{tab:surrogate-point}
\begin{tabular}{@{}lrrr@{}}
\toprule
Target & MAE & RMSE & $R^2$ \\
\midrule
HOSR           & 0.043 & 0.075 & 0.956 \\
RLF rate       & 0.004 & 0.005 & 0.952 \\
Ping-pong rate & 0.015 & 0.019 & 0.717 \\
\bottomrule
\end{tabular}
\end{table}

\paragraph{Uncertainty calibration.}
Point accuracy is not enough for a recommendation; the intervals must be
trustworthy. Table~\ref{tab:surrogate-cov} shows that the naive quantile
intervals are too narrow under grouped cross-validation: their empirical coverage
of the nominal $90\%$ interval falls to $0.65$--$0.83$. Split-conformal
calibration moves all three targets back towards nominal ($0.72$--$0.89$), at the
cost of slightly wider intervals. Figure~\ref{fig:reliability} shows the same
correction as a reliability diagram. This is why the end-to-end loop reports
conformal rather than naive intervals.

\begin{table}[ht]
\centering
\caption{Empirical coverage of the nominal 90\% prediction interval (own elaboration).}
\label{tab:surrogate-cov}
\begin{tabular}{@{}lrr@{}}
\toprule
Target & Naive coverage & Conformal coverage \\
\midrule
HOSR           & 0.77 & 0.84 \\
RLF rate       & 0.65 & 0.72 \\
Ping-pong rate & 0.83 & 0.89 \\
\bottomrule
\end{tabular}
\end{table}

\begin{figure}[ht]
\centering\includegraphics[width=.85\textwidth]{reliability_diagram}
\caption{Naive vs split-conformal interval coverage (own elaboration).}
\label{fig:reliability}
\end{figure}

\paragraph{Inverse query.}
Asked which configuration maximises HOSR under the operating constraints
(HOSR~$\geq0.95$, RLF~$\leq0.05$, ping-pong~$\leq0.10$) at the median deployment
context, the surrogate finds $12$ of the $36$ candidates feasible. The top three
all use the shortest time-to-trigger, TTT~$=256$\,ms (with hysteresis/A3 offset
$(6,0)$, $(2,3)$ and $(0,6)$\,dB), each reaching predicted HOSR~$\approx 1.0$.
The interpretation is that, in this radio context, a fast trigger paired with
enough hysteresis or offset to suppress ping-pong is preferable, and the
surrogate makes that trade-off explicit (Figure~\ref{fig:inverse}).

\begin{figure}[ht]
\centering\includegraphics[width=.85\textwidth]{inverse_query}
\caption{Feasible configurations on the HOSR vs RLF plane; top-3 annotated (own elaboration).}
\label{fig:inverse}
\end{figure}

\paragraph{End-to-end reconfiguration.}
Running the full loop over the timeline produces seven drift-triggered
interventions. The most relevant case is the harmful reconfiguration drift
(D-4), where the operator raises the time-to-trigger: the monitor detects it and
the surrogate recommends a recovery configuration, and across all interventions
the cumulative \emph{predicted} HOSR uplift is $+0.62$
(Figure~\ref{fig:end2end}). All four demo acceptance checks pass: at least five
interventions fire, every recommendation lies inside the trained grid, both D-4
events change the configuration, and the cumulative HOSR uplift is positive. This
completes the RQ4 loop: drift can degrade the handover KPI, and a feasible
reconfiguration, chosen from the surrogate, is predicted to recover it. The
uplift is a model prediction, not a measured field gain, which is stated
as a limitation.

\begin{figure}[ht]
\centering\includegraphics[width=.95\textwidth]{end_to_end_moneyshot}
\caption{End-to-end loop: sliding PR-AUC, anomaly score, interventions and cumulative HOSR uplift (own elaboration).}
\label{fig:end2end}
\end{figure}

\section{External validation}
\label{sec:res-external}

How far do these results travel? Table~\ref{tab:external} summarises the three
checks. First, a cross-seed replication over five master
seeds: the drift-volume is stable (coefficient of variation $0.04$), the
filtered-beats-naive gap replicates ($+0.072$, 95\% CI $[+0.033, +0.113]$), the
end-to-end HOSR uplift is essentially identical across seeds
($+0.621$), and the pooled anomaly winner is Isolation Forest in all five seeds.
The one criterion that does not pass is the stability of the \emph{per-type}
anomaly winner, which alternates between LOF, One-Class SVM and MLP-AE across
seeds; this is consistent with the tight per-type cluster seen in
Section~\ref{sec:res-anomaly} and is reported transparently as a limitation.

Second, a cross-scenario check on a denser urban timeline: all five criteria
pass. The filtered-beats-naive sign is preserved, and the end-to-end uplift not
only holds but scales to $+1.24$, consistent with there being more to gain when
mobility is more demanding. Third, a face-validity check on real NordicDat
5G-NSA telemetry, run fully self-supervised: the monitor flags only $1.2\%$ of
windows (well under the $10\%$ sanity ceiling) and its flagging peaks at the
$13{:}00$ commute hour ($3.5\times$ the daily mean), which is operationally
plausible. Because that data has no ground truth, no precision or recall is
claimed.

\begin{table}[ht]
\centering
\caption{External-validation headline outcomes (own elaboration).}
\label{tab:external}
\begin{tabular}{@{}lll@{}}
\toprule
Check & Result & Verdict \\
\midrule
Cross-seed: filtered $>$ naive      & $+0.072$ [$+0.033,+0.113$] & pass \\
Cross-seed: end-to-end HOSR uplift  & $+0.621$ (CV $\approx 0$)  & pass \\
Cross-seed: per-type anomaly winner & unstable (LOF/OCSVM/MLP-AE) & fail \\
Cross-scenario (dense urban)        & all 5 criteria; uplift $+1.24$ & pass \\
NordicDat face validity             & $1.2\%$ flagged, $13{:}00$ peak & pass \\
\bottomrule
\end{tabular}
\end{table}

\section{Discussion}
\label{sec:res-discussion}

Taken together, the four layers are consistent with the thesis premise that
handover monitoring benefits from being anomaly- and drift-aware rather than
trained once and left unchanged. Anomaly detection works well for structured
faults but struggles with subtle, single-UE or short cell-local ones, which sets
a realistic expectation for what an unsupervised monitor can catch. Drift
detection is dependable but involves an explicit latency-versus-false-alarm
choice. The adaptive layer suggests that retraining after drift can make things
worse if applied naively, and that a simple self-supervised filter is what allows
retraining to yield a net gain. Finally, the surrogate links configuration to
outcome accurately enough to suggest a recovery configuration after a harmful
change, with calibrated uncertainty. The main
limitations follow directly from the design: the anomalies are injected rather
than observed in the field, the KPI uplift is a surrogate prediction rather than
a measured gain, and the per-type anomaly ranking is seed-sensitive. These are
revisited as future work in the next chapter.
